\subsection{复测度的绝对值}

设 \(\mu\) 是局部紧空间 \(X\) 上的一个复测度；对每个函数 \(f \in \mathcal{K}_{+}(X)\) ，正实数

\[
L (f) = \sup _ {| g | \leqslant f, g \in \mathcal {K} (\mathrm{X}; \mathbf {C})} \left| \int g d \mu \right|\tag{11}
\]

是有限的，因为关系 \(|g| \leqslant f\) 蕴涵 \(\operatorname{Supp}(g) \subset \operatorname{Supp}(f)\) 且 \(\| g \| \leqslant \| f \|\) ，于是我们的断言由 No. 3 的公式 (4) 得出。我们证明 \(L\) 可以以唯一方式延拓成 \(X\) 上的一个正测度；鉴于 No. 5 的定理 1 与 Ch. II, §2, No. 1, 命题 3，只须证明：若 \(f_1, f_2\) 是 \(\mathcal{K}_+(X)\) 中的两个函数，则 \(L(f_1 + f_2) = L(f_1) + L(f_2)\) 。现在，若 \(|g_1| \leqslant f_1\) 且 \(|g_2| \leqslant f_2\) ，其中 \(g_1\) 与 \(g_2\) 是 \(\mathcal{K}(X; \mathbf{C})\) 中的函数，则有 \(|g_1 + \zeta g_2| \leqslant f_1 + f_2\) （对任何绝对值为 1 的复数 \(\zeta\) ），因此

\[
\left| \mu (g _ {1} + \zeta g _ {2}) \right| = \left| \mu (g _ {1}) + \zeta \mu (g _ {2}) \right| \leqslant L (f _ {1} + f _ {2}).
\]

此外，可以假定 \(\zeta\) 如此选取，使得

\[
\left| \mu (g _ {1}) + \zeta \mu (g _ {2}) \right| = \left| \mu (g _ {1}) \right| + \left| \mu (g _ {2}) \right|;
\]

由于 \(|\mu(g_i)|\) 可任意接近 \(L(f_i)\) （ \(i = 1,2\) ），这就证明了 \(L(f_1) + L(f_2) \leqslant L(f_1 + f_2)\) 。另一方面，考虑一个函数 \(g \in \mathcal{K}(\mathbf{X};\mathbf{C})\) ，它满足 \(|g| \leqslant f_1 + f_2\) 。函数 \(g_i\) ——它等于 \(g f_i / (f_1 + f_2)\) 于使 \(f_1(x) + f_2(x) \neq 0\) 的点处，在别处等于 0（ \(i = 1,2\) ）——在 \(\mathbf{X}\) 上连续，因为 \(f_i / (f_1 + f_2)\) （ \(i = 1,2\) ）在使 \(f_1(x) + f_2(x) \neq 0\) 的每一点处连续，且 \(|g_i(x)| \leqslant |g(x)|\) 对每个 \(x \in \mathbf{X}\) 成立，这证明了 \(g_i\) 在使 \(f_1(x) + f_2(x) = 0\) 的点处的连续性（ \(i = 1,2\) ），因为在这些点处也有 \(g(x) = 0\) 。显然 \(|g_i| \leqslant f_i\) （ \(i = 1,2\) ）且 \(g = g_1 + g_2\) ，因此

\[
| \mu (g) | \leqslant | \mu (g _ {1}) | + | \mu (g _ {2}) | \leqslant L (f _ {1}) + L (f _ {2});
\]

由于 \(|\mu(g)|\) 可任意接近 \(L(f_1 + f_2)\) ，我们有

\[
L (f _ {1} + f _ {2}) \leqslant L (f _ {1}) + L (f _ {2}),
\]

这就完成了我们断言的证明。

当 \(\mu\) 是实测度时，由公式 (9) 可知 \(|\mu| \leqslant L\) ；另一方面，由 No. 5 的末尾部分，若 \(g \in \mathcal{K}(\mathbf{X}; \mathbf{C})\) 且 \(|g| \leqslant f \in \mathcal{K}_{+}(\mathbf{X})\) ，则 \(|\int g d\mu| \leqslant \int |g| \cdot d|\mu| \leqslant \int f d|\mu|\) ，因此按定义 \(L \leqslant |\mu|\) ，换言之 \(L = |\mu|\) 。

我们仍用 \(|\mu|\) 表示正测度 L（对任何复测度 \(\mu\) ），并称 \(|\mu|\) 为 \(\mu\) 的绝对值。于是 \(|\mu|\) 的定义可以写成

\[
| \mu | (f) = \sup _ {| g | \leqslant f, g \in \mathcal{K} (\mathrm{X}; \mathbf {C})} | \mu (g) |,\tag{12}
\]

因而，对每个函数 \(g \in \mathcal{K}(\mathrm{X};\mathbf{C})\)

\[
\left| \int g d \mu \right| \leqslant \int | g | d | \mu |.\tag{13}
\]

显然，对每个标量 \(\alpha \in \mathbf{C}\) 及每个测度 \(\mu\) （\(X\) 上的），

\[
| \alpha \mu | = | \alpha | \cdot | \mu |.\tag{14}
\]

另一方面，若 \(\mu\) 与 \(\nu\) 是 \(X\) 上的两个测度，\(f\) 是 \(\mathcal{K}_{+}(X)\) 中的一个函数，且 \(g\) 是 \(\mathcal{K}(X; \mathbf{C})\) 中满足 \(|g| \leqslant f\) 的一个函数，则

\[
\left| \int g d (\mu + \nu) \right| = \left| \int g d \mu + \int g d \nu \right| \leqslant \int f d | \mu | + \int f d | \nu |,
\]

由此

\[
| \mu + \nu | \leqslant | \mu | + | \nu |.\tag{15}
\]

用同样的记号，关系 \(|g| \leqslant f\) 与 \(|\overline{g}| \leqslant f\) 等价，因此

\[
| \overline {{{\mu}}} | = | \mu |.\tag{16}
\]

由 (14)、(15) 与 (16) 推出

\[
| \mathcal {R} \mu | \leqslant | \mu |, \quad | \mathcal {I} \mu | \leqslant | \mu |, \quad | \mu | \leqslant | \mathcal {R} \mu | + | \mathcal {I} \mu |.\tag{17}
\]

命题 8. --- 若 \(\mu\) 是 X 上的一个测度，则对每个函数 \(h \in \mathcal{C}(\mathrm{X};\mathbf{C})\) ，

\[
\left| h \cdot \mu \right| \leqslant \left| h \right| \cdot \left| \mu \right|.\tag{18}
\]

因为，若 \(f \in \mathcal{K}_+(\mathbf{X})\) 且 \(g \in \mathcal{K}(\mathbf{X};\mathbf{C})\) 满足 \(|g| \leqslant f\) ，则由 (13)，\(|\int gh d\mu| \leqslant \int |gh| d|\mu| \leqslant \int f|h| d|\mu|\) ，这就证明了 (18)。

